\chapter{Glossário inglês--português}
\label{ap:glossario}

Os termos em inglês explicados em nota de rodapé ao longo da apostila (só
eles), em ordem alfabética do inglês. Quando o texto usa o termo em
português, ele está na segunda coluna.

\small
\begin{longtable}{@{}>{\raggedright}p{3.9cm} >{\raggedright}p{3.6cm} >{\raggedright\arraybackslash}p{6.3cm}@{}}
\toprule
Inglês & Português & Em uma linha \\
\midrule \endhead
\textit{accessors} & métodos de acesso & Métodos de leitura; num record, \texttt{x()} e não \texttt{getX()} \\
\textit{adapter} & adaptador & Converte a API da porta para um dispositivo concreto (HTTP, SQL, fila) \\
\textit{augmentation} & aumento & Fase do build do Quarkus em que as extensões geram o código de runtime \\
\textit{BEAM} & (nome próprio) & A máquina virtual do Erlang \\
\textit{blast radius} & raio de explosão & O tanto do sistema e dos usuários que um experimento pode afetar \\
\textit{bootstrap} & método de inicialização & Chamado uma vez, na primeira execução de um \texttt{invokedynamic} \\
\textit{build items} & itens do build & Objetos que um passo do build produz e outros consomem \\
\textit{build steps} & passos do build & Métodos \texttt{@BuildStep} de uma extensão do Quarkus \\
\textit{class loader} & carregador de classes & Encontra os bytes de uma classe e a define na JVM \\
\textit{CLI (Common Language Infrastructure)} & infraestrutura comum de linguagem & Padrão de VM multilinguagem da plataforma .NET \\
\textit{client proxy} & proxy do cliente & Objeto gerado no lugar do bean injetado, que delega à instância real \\
\textit{closed-world assumption} & suposição de mundo fechado & Todo o código que pode executar é conhecido no build \\
\textit{code smells} & cheiros de código & Sinais de código mal estruturado, mesmo que funcione \\
\textit{constant pool} & tabela de constantes & Nomes, strings e números de um \texttt{.class} \\
\textit{coordinated omission} & omissão coordenada & O gerador de carga espera a resposta lenta e deixa de medir a espera real \\
\textit{data race} & corrida de dados & Acessos concorrentes, ao menos um de escrita, sem happens-before \\
\textit{deoptimization} & desotimização & Trocar, em execução, código otimizado por código não otimizado \\
\textit{escape analysis} & análise de escape & Descobre se um objeto sai do método ou da thread que o criou \\
\textit{event loop} & laço de eventos & Thread que atende muitas conexões sem bloquear \\
\textit{fixture} & (sem tradução usual) & Estado inicial conhecido de que o teste precisa \\
\textit{garbage collector (GC)} & coletor de lixo & Recupera a memória de objetos que não são mais alcançáveis \\
\textit{Gizmo} & (nome próprio) & Biblioteca de geração de bytecode do Quarkus \\
\textit{happens-before} & acontece antes & Relação que garante visibilidade e ordem entre ações de threads \\
\textit{hot code loading} & carga de código a quente & Trocar um módulo com o sistema rodando \\
\textit{inline cache} & cache em linha & Resultado da busca de método guardado no ponto de chamada \\
\textit{inlining} & expansão em linha & Copiar o corpo do método chamado para dentro do chamador \\
\textit{Jandex} & (nome próprio) & Índice de arquivos \texttt{.class} da SmallRye \\
\textit{JIT (just-in-time)} & compilação durante a execução & Compila com informação de como o programa está rodando \\
\textit{let it crash} & deixe quebrar & O processo termina e outro cuida da recuperação \\
\textit{load shedding} & descarte de carga & Descartar requisições menos importantes sob sobrecarga \\
\textit{mailbox} & caixa de correio & Fila de mensagens de um processo Erlang \\
\textit{mock} & (sem tradução usual) & Dublê que verifica as chamadas que recebe \\
\textit{p-code (pseudo-code)} & pseudocódigo & Instruções da P-machine do Pascal, executadas por interpretador \\
\textit{partial evaluation} & avaliação parcial & Especializar um programa para entradas já conhecidas \\
\textit{PGO (profile-guided optimization)} & otimização guiada por perfil & AOT guiado por um perfil de uma execução de treino \\
\textit{pinned} & presa & Virtual thread que não consegue desmontar da carrier thread \\
\textit{polymorphic inline caches (PICs)} & caches em linha polimórficos & Inline cache com uma entrada por tipo de receptor \\
\textit{port} & porta & Ponto em que a aplicação conversa com o mundo de fora \\
\textit{preview} & prévia & Recurso para experimentação, com \texttt{-{}-enable-preview} \\
\textit{recorders} & gravadores & Chamadas gravadas no build e reproduzidas na subida da aplicação \\
\textit{reductions} & reduções & Unidades de trabalho de um processo Erlang \\
\textit{RSS (resident set size)} & memória residente & Memória física de fato ocupada pelo processo \\
\textit{safepoints} & pontos seguros & Pontos em que a JVM conhece o estado da thread e pode pará-la \\
\textit{schedulers} & escalonadores & Threads do SO que executam processos Erlang \\
\textit{stub} & (sem tradução usual) & Dublê que devolve respostas prontas \\
\textit{supervision tree} & árvore de supervisão & Supervisores e seus processos filhos \\
\textit{SUT (system under test)} & sistema sob teste & O código que o teste está exercitando \\
\textit{throughput} & vazão & Requisições atendidas por unidade de tempo \\
\textit{tiered compilation} & compilação em camadas & Interpretador, depois C1 com perfil, depois C2 \\
\textit{uncommon trap} & (sem tradução usual) & Devolve o controle ao runtime num caminho tido como impossível \\
\textit{unmount} & desmontar & A virtual thread sai da carrier thread, que fica livre \\
\bottomrule
\end{longtable}
\normalsize
